% ---------------------------------------------------------------------------
% Appendix: proofs for Sec. \ref{subsec:idm-theory}. Rewritten 2026-11-04
% (author: "remove the unnecessary stuff ... make the proofs more clear"):
% kept only what Sec. 4.4 uses; proofs in numbered steps. Removed (see git
% history, or scratchpad/appendix_idm_before_cleanup.tex of that session):
% the 2x2 counterexample to a Euclidean contraction, the noise-instantiation
% bound, the covariate-shift density bound, the Gaussian/sub-Gaussian tail
% bound (old Prop. (c)), the general one-sided covering, the union bound over
% a trial, and the tightness part (d) of the grid lemma with its numerical
% search. Numerical checks of the original: scratchpad/idmcheck/*.py.
% ---------------------------------------------------------------------------
\section{Analysis of the Inverse Footstep Model}
\label{app:idm-proofs}

\looseness=-1 This appendix proves the statements of
Sec.~\ref{subsec:idm-theory}: when the regression of
Eq.~\eqref{eq:idm-loss} inverts the controller (Proposition~\ref{prop:idm-mean}),
how it is biased otherwise (Proposition~\ref{prop:idm-shrink}), when a
landing error can put the foot in a void (Proposition~\ref{prop:idm-void}),
and how much clearance the grown mask leaves (Lemma~\ref{lem:idm-grid}).
We drop the superscript $c$ from $I^c$. Let $P_c$ be the distribution of the
logged tuples $(\zeta, \mathbf u, \mathbf o)$, one per step the controller
executed, and $I_\star(\zeta,\mathbf o) := \mathbb E_{P_c}[\mathbf u \mid
\zeta,\mathbf o]$ the function the regression estimates. Norms are
Euclidean unless marked, $\mathrm{Id}$ is the identity matrix, and $\Phi$
is the standard normal distribution function. We assume $\mathbb
E_{P_c}\|\mathbf u\|^2 < \infty$ throughout.

\looseness=-1 Throughout, we assume that the law of the landing $\mathbf o$
given $(\zeta,\mathbf u)$ is the same when logging and at deployment,
although the logging policy and GaitNet both choose commands from terrain
that $\zeta$ omits. Two parts of our pipeline strain our assumptions. The
training controller's IDM is fitted without the controller randomization
it later runs under (Sec.~\ref{subsec:method-training}), which changes
that law. The OCS2 IDM is fitted on steps sent through the training
controller's IDM, so its logged landings need not cover the deployed
targets, as Assumption~\ref{as:idm-exact}(iii) below requires; logging
again with the deployed IDM \citep{christiano_transfer_2016} would reduce
this gap.

\subsection{When the Regression Inverts the Controller}

\begin{assumption}[Exact relabelling]
\label{as:idm-exact}
(i)~There is no residual: the landing is $\mathbf o = \bar
F_c(\zeta,\mathbf u)$ for every $(\zeta,\mathbf u)$ in the support of the
logged pairs, when logging and at deployment.
(ii)~For every $(\zeta,\mathbf o)$, the commands $\mathbf u$ with
$(\zeta,\mathbf u)$ in that support and $\bar F_c(\zeta,\mathbf u) =
\mathbf o$ form a convex set.
(iii)~Every deployment pair $(\zeta,\mathbf u^\star)$ lies in the support
of the logged $(\zeta,\mathbf o)$, and $\bar F_c$ (on the support of the
logged pairs) and $I_\star$ have continuous versions, which we use.
(iv)~The fitted IDM $I$ is continuous and equals $I_\star$ $P_c$-almost
everywhere.
\end{assumption}

Condition (ii) holds, with single points, whenever the landing offset
changes less than the command, that is, when $\mathbf u \mapsto \bar
F_c(\zeta,\mathbf u) - \mathbf u$ is Lipschitz with a constant below one
(step~5 of the proof below).

\begin{proposition}[Hindsight regression]
\label{prop:idm-mean}
(a)~Among the square-integrable functions of $(\zeta,\mathbf o)$,
$I_\star$ minimises the loss of Eq.~\eqref{eq:idm-loss}, uniquely up to
$P_c$-null sets.
(b)~Under Assumption~\ref{as:idm-exact}, every deployment pair has
$\bar F_c(\zeta, I(\zeta,\mathbf u^\star))$ $= \mathbf u^\star$, so the foot
lands at $\mathbf u^\star$.
\end{proposition}

Every logged step is an example of the inverse map, since its command
landed the foot at $\mathbf o$; noise added to a command before it is
sent therefore does not corrupt its label. Each condition of
Assumption~\ref{as:idm-exact} matters. Without (i), the regression averages
$\mathbf u$ over the hidden states that share $(\zeta,\mathbf o)$ and is
biased even with unlimited data \citep{eysenbach_imitating_2022}
(Proposition~\ref{prop:idm-shrink}). Without (ii), commands from two
inverse branches that land at the same point can average to one that
lands at neither \citep{jordan_forward_1992}. Without (iii), the held-out
error says nothing about targets outside the logged landings.

\begin{idmproof}{Proposition~\ref{prop:idm-mean}}
(a)~For a square-integrable $I$, split $\mathbf u - I = (\mathbf u -
I_\star) + (I_\star - I)$. The cross term vanishes, $\mathbb E[(\mathbf u -
I_\star)^\top(I_\star - I)] = \mathbb E\big[\mathbb E[\mathbf u - I_\star
\mid \zeta,\mathbf o]^\top (I_\star - I)\big] = 0$, because $I_\star - I$
depends only on $(\zeta,\mathbf o)$ and $\mathbb E[\mathbf u - I_\star
\mid \zeta,\mathbf o] = \mathbf 0$. Hence
\begin{equation}
  \mathbb E\|\mathbf u - I\|^2 = \mathbb E\|\mathbf u - I_\star\|^2
  + \mathbb E\|I_\star - I\|^2 ,
  \label{eq:app-orth}
\end{equation}
which is smallest exactly when $I = I_\star$ almost everywhere. In words,
the held-out squared error of any IDM is the irreducible error plus its
mean squared distance from $I_\star$.

(b)~We use that the support of a distribution is the smallest closed
set of probability one, so a closed set of probability one contains it.

\emph{Step 1: the logged commands for a landing.} Let $\mathcal
U(\zeta,\mathbf o)$ be the set of commands $\mathbf u$ with $(\zeta,\mathbf
u)$ in the support of the logged pairs and $\bar F_c(\zeta,\mathbf u) =
\mathbf o$. It is closed because the support is closed and $\bar F_c$ is
continuous, and convex by (ii). Logged pairs lie in their support and, by
(i), land at $\bar F_c(\zeta,\mathbf u)$, so given $(\zeta,\mathbf o)$ the
command lies in $\mathcal U(\zeta,\mathbf o)$ with probability one.

\emph{Step 2: the mean stays in the set.} The mean of a distribution that
gives probability one to a closed convex set lies in that set: if the
mean $\bar{\mathbf x}$ were outside, a separating hyperplane would give a vector
$\boldsymbol\nu$ and a number $t$ with $\boldsymbol\nu^\top\bar{\mathbf x} > t
\ge \boldsymbol\nu^\top\mathbf x$ for every $\mathbf x$ in the set, and
taking the mean of the second inequality would give
$\boldsymbol\nu^\top\bar{\mathbf x} \le t$. Applied to the conditional
distribution of $\mathbf u$ given $(\zeta,\mathbf o)$, this gives
$I_\star(\zeta,\mathbf o) \in \mathcal U(\zeta,\mathbf o)$, so $\bar
F_c(\zeta, I_\star(\zeta,\mathbf o)) = \mathbf o$ for $P_c$-almost every
$(\zeta,\mathbf o)$.

\emph{Step 3: from almost every landing to every target.} By step~2,
the set of $(\zeta,\mathbf o)$ with $(\zeta, I_\star(\zeta,\mathbf o))$ in
the support of the logged pairs and $\bar F_c(\zeta, I_\star(\zeta,\mathbf
o)) = \mathbf o$ has probability one. It is closed, as $I_\star$ and $\bar
F_c$ are continuous by (iii), so it contains the support of the logged
$(\zeta,\mathbf o)$, and with it every deployment pair $(\zeta,\mathbf
u^\star)$. Hence the command $I_\star(\zeta,\mathbf u^\star)$ lies in the
support of the logged pairs and $\bar F_c(\zeta, I_\star(\zeta,\mathbf
u^\star)) = \mathbf u^\star$.

\emph{Step 4: the fitted IDM.} By (iv), the set where the continuous $I$
equals $I_\star$ is closed and has probability one, so it contains the
support, and $I(\zeta,\mathbf u^\star) = I_\star(\zeta,\mathbf u^\star)$.
This command lies in the support of the logged pairs, where by (i) the
landing equals $\bar F_c$, so the foot lands at $\mathbf u^\star$.

\emph{Step 5: the condition after the assumption.} If $\mathbf u \mapsto
\bar F_c(\zeta,\mathbf u) - \mathbf u$ is $L_\delta$-Lipschitz with
$L_\delta < 1$, then $\|\bar F_c(\zeta,\mathbf u) - \bar F_c(\zeta,\mathbf
u')\| \ge \|\mathbf u - \mathbf u'\| - L_\delta\|\mathbf u - \mathbf u'\| >
0$ for $\mathbf u \ne \mathbf u'$, so $\bar F_c(\zeta,\cdot)$ is injective
and every $\mathcal U(\zeta,\mathbf o)$ has at most one point.
\end{idmproof}

\subsection{Bias When the Landings Are Noisy}

\begin{proposition}[Relabelling bias]
\label{prop:idm-shrink}
Fix $\zeta$. Let the logged commands be $\mathbf u \sim \mathcal
N(\boldsymbol\mu,\Sigma_u)$ with $\Sigma_u \succ 0$, and the landings
$\mathbf o = A\mathbf u + \mathbf b + \boldsymbol\eta$ with $A$ invertible
and $\boldsymbol\eta \sim \mathcal N(\mathbf 0,\Sigma_\eta)$ independent
of $\mathbf u$. Let $\bar{\mathbf o} = A\boldsymbol\mu + \mathbf b$ be the
mean logged landing and $S = A\Sigma_u A^\top + \Sigma_\eta$ the covariance
of the logged landings. For a target $\mathbf u^\star$, the landing
$\mathbf o^\star$ of the command $I_\star(\zeta,\mathbf u^\star)$ has the
mean error
\begin{equation}
  \mathbb E[\mathbf o^\star] - \mathbf u^\star
  = -\Lambda\,(\mathbf u^\star - \bar{\mathbf o}), \qquad
  \Lambda = \Sigma_\eta S^{-1}.
  \label{eq:idm-shrink}
\end{equation}
The eigenvalues of $\Lambda$ lie in $[0,1)$, and
$\|\mathbb E[\mathbf o^\star] - \bar{\mathbf o}\|_{S^{-1}} \le \|\mathbf
u^\star - \bar{\mathbf o}\|_{S^{-1}}$, with $\|\mathbf x\|_{S^{-1}}^2 :=
\mathbf x^\top S^{-1}\mathbf x$. A larger $\Sigma_u$ (in the
positive-semidefinite order) raises no eigenvalue of $\Lambda$.
\end{proposition}

\looseness=-1 Here $\Sigma_\eta$ is the part of the landing covariance $S$
that the command does not explain. In one dimension, $\Lambda$ is that
share of the landing variance, and the deployed landing falls short of
$\mathbf u^\star$, toward the mean logged landing $\bar{\mathbf o}$, by
the fraction $\Lambda$ of their distance. In the plane, the pull toward
$\bar{\mathbf o}$ holds in the metric $\|\cdot\|_{S^{-1}}$ of the logged
landings, not necessarily in the Euclidean one, and the bias, measured
in that metric, is at most the largest eigenvalue of $\Lambda$ times
$\|\mathbf u^\star - \bar{\mathbf o}\|_{S^{-1}}$, so it can be largest for footholds that the logging policy
seldom reached. This is the attenuation bias of a regression on a noisy
input. Spreading the logged commands raises no eigenvalue of $\Lambda$,
and in one dimension it shrinks the bias, a second reason, besides
coverage, for the command noise in the training controller's IDM data
(App.~\ref{app:settings});
in the plane, the bias at a given target can still grow.

\begin{idmproof}{Proposition~\ref{prop:idm-shrink}}
\emph{Step 1: the regression is linear.} $(\mathbf u,\mathbf o)$ is
jointly Gaussian with $\operatorname{Cov}(\mathbf o) = S \succ 0$ and
$\operatorname{Cov}(\mathbf u,\mathbf o) = \Sigma_u A^\top$, so
$I_\star(\zeta,\mathbf o) = \boldsymbol\mu + W(\mathbf o - \bar{\mathbf
o})$ with the gain $W = \Sigma_u A^\top S^{-1}$.

\emph{Step 2: the mean landing.} The command $\tilde{\mathbf u} =
I_\star(\zeta,\mathbf u^\star)$ is fixed at the decision and
$\boldsymbol\eta$ is independent of it, so $\mathbb E[\mathbf o^\star] =
A\tilde{\mathbf u} + \mathbf b = \bar{\mathbf o} + AW(\mathbf u^\star -
\bar{\mathbf o})$. Hence $\mathbb E[\mathbf o^\star] - \mathbf u^\star =
-(\mathrm{Id} - AW)(\mathbf u^\star - \bar{\mathbf o})$, and $\mathrm{Id} -
AW = (S - A\Sigma_u A^\top)S^{-1} = \Sigma_\eta S^{-1} = \Lambda$.

\emph{Step 3: the eigenvalues.} Let $M := S^{-1/2}\Sigma_\eta S^{-1/2} =
\mathrm{Id} - S^{-1/2}A\Sigma_u A^\top S^{-1/2}$. It is symmetric, positive
semidefinite because $\Sigma_\eta \succeq 0$, and $M \prec \mathrm{Id}$ because
$A\Sigma_u A^\top \succ 0$, so its eigenvalues lie in $[0,1)$. Since
$\Lambda = S^{1/2} M S^{-1/2}$, $\Lambda$ has the same eigenvalues.

\emph{Step 4: the contraction.} $\mathbb E[\mathbf o^\star] - \bar{\mathbf
o} = (\mathrm{Id} - \Lambda)(\mathbf u^\star - \bar{\mathbf o})$ and
$S^{-1/2}(\mathrm{Id} - \Lambda) = (\mathrm{Id} - M)S^{-1/2}$. As $0 \prec
\mathrm{Id} - M \preceq \mathrm{Id}$, $\|\mathbb E[\mathbf o^\star] -
\bar{\mathbf o}\|_{S^{-1}} = \|(\mathrm{Id} - M)S^{-1/2}(\mathbf u^\star -
\bar{\mathbf o})\| \le \|S^{-1/2}(\mathbf
u^\star - \bar{\mathbf o})\| = \|\mathbf u^\star - \bar{\mathbf
o}\|_{S^{-1}}$. In the same way, $\|S^{-1/2}\Lambda(\mathbf u^\star -
\bar{\mathbf o})\| = \|M S^{-1/2}(\mathbf u^\star - \bar{\mathbf o})\| \le
\lambda_{\max}(M)\,\|\mathbf u^\star - \bar{\mathbf o}\|_{S^{-1}}$, which
bounds the bias in the metric $\|\cdot\|_{S^{-1}}$.

\emph{Step 5: more command spread.} The eigenvalues of $\Lambda =
\Sigma_\eta S^{-1}$ equal those of the symmetric matrix
$\Sigma_\eta^{1/2}S^{-1}\Sigma_\eta^{1/2}$, since $XY$ and $YX$ have the
same eigenvalues. A larger $\Sigma_u$ gives a larger $S$ and so a smaller
$S^{-1}$ and a smaller $\Sigma_\eta^{1/2}S^{-1}\Sigma_\eta^{1/2}$, whose
ordered eigenvalues therefore do not increase.
\end{idmproof}

\subsection{Landing Error and Voids}

\begin{proposition}[Landing error and voids]
\label{prop:idm-void}
For the command $\mathbf u = I(\zeta,\mathbf u^\star)$ and the clearance
$d(\mathbf u^\star)$ of Sec.~\ref{subsec:idm-theory}, let $\mathcal G$ be
the information at the decision, which fixes $\zeta$, $\mathbf u^\star$,
$\mathbf u$ and $\mathcal V$, and let $\boldsymbol\beta := \bar
F_c(\zeta,\mathbf u) - \mathbf u^\star$. Then the landing error is
$\mathbf e := \mathbf o - \mathbf u^\star = \boldsymbol\beta +
\boldsymbol\eta$, and
(a)~if $\|\boldsymbol\beta\| + \|\boldsymbol\eta\| < d(\mathbf u^\star)$,
the foot does not land in a void;
(b)~$\Pr(\mathbf o \in \mathcal V \mid \mathcal G) \le \Pr(\|\mathbf
e\|_\infty \ge d(\mathbf u^\star) \mid \mathcal G)$.
\end{proposition}

\begin{idmproof}{Proposition~\ref{prop:idm-void}}
By Eq.~\eqref{eq:mean-forward}, $\mathbf e = \bar F_c(\zeta,\mathbf u) +
\boldsymbol\eta - \mathbf u^\star = \boldsymbol\beta + \boldsymbol\eta$. As $\boldsymbol\eta$ has mean zero given
$(\zeta,\mathbf u)$, the mean error is $\boldsymbol\beta$, so an IDM with
the inputs $\zeta$ can cancel at most $\boldsymbol\beta$.
If $\mathbf o \in \mathcal V$, then, since $d(\mathbf u^\star)$ is the
$\infty$-norm distance from $\mathbf u^\star$ to the nearest void point,
$d(\mathbf u^\star) \le \|\mathbf o - \mathbf u^\star\|_\infty = \|\mathbf
e\|_\infty \le \|\mathbf e\| \le \|\boldsymbol\beta\| +
\|\boldsymbol\eta\|$. Read backwards, this is (a). It also shows $\{\mathbf
o \in \mathcal V\} \subseteq \{\|\mathbf e\|_\infty \ge d(\mathbf
u^\star)\}$, which is (b). Both statements weaken as $d$ shrinks, so a
lower bound on $d(\mathbf u^\star)$, such as that of
Lemma~\ref{lem:idm-grid}, may replace it.
\end{idmproof}

\paragraph{Where $\boldsymbol\beta$ comes from.}
\looseness=-1 With $\tilde{\mathbf u} := I_\star(\zeta,\mathbf u^\star)$,
\begin{equation}
  \boldsymbol\beta = \big[\bar F_c(\zeta,\mathbf u) - \bar
  F_c(\zeta,\tilde{\mathbf u})\big] + \big[\bar F_c(\zeta,\tilde{\mathbf
  u}) - \mathbf u^\star\big].
  \label{eq:app-beta}
\end{equation}
If $\bar F_c(\zeta,\cdot)$ is $L$-Lipschitz on the segment from
$\tilde{\mathbf u}$ to $\mathbf u$, the first term is at most
$L\|I(\zeta,\mathbf u^\star) - I_\star(\zeta,\mathbf u^\star)\|$: the
IDM's fitting error, whose mean square over the logged landings is its
held-out error minus the irreducible error (Eq.~\eqref{eq:app-orth}); at
a given target, the held-out error speaks for it only where the logged
landings cover that target (Assumption~\ref{as:idm-exact}(iii)). In the form of
Eq.~\eqref{eq:forward}, $\bar F_c(\zeta,\mathbf u) = \mathbf u +
\mathbb E[\boldsymbol\delta_c \mid \zeta,\mathbf u]$, so a landing offset
that depends only weakly on the command gives $L \approx 1$. The second
term, the relabelling bias, vanishes under Assumption~\ref{as:idm-exact}
(steps 1--3 of the proof of Proposition~\ref{prop:idm-mean}) and equals
$-\Lambda(\mathbf u^\star - \bar{\mathbf o})$ in the setting of
Proposition~\ref{prop:idm-shrink}.

\begin{remark}[Landing error against the margin]
\label{rem:idm-numbers}
On flat ground, the IDM cut OCS2's landing error from
\corlOcsLandRmsNoIdm\ to \corlOcsLandRmsIdm\,cm RMS
(Sec.~\ref{subsec:exp-why}). Suppose $\mathbf e$ is zero-mean isotropic
Gaussian with standard deviation $\sigma$ per axis. A foothold whose
nearest void lies beyond a long straight edge at Euclidean distance
$c_\perp$ then lands in it with probability $\Phi(-c_\perp/\sigma)$,
because only the error component normal to the edge matters. An edge at
$\infty$-norm distance $d$ with unit normal $\mathbf n$ has $c_\perp =
d\|\mathbf n\|_1 \ge d$. For our candidates, Lemma~\ref{lem:idm-grid}(c)
with $k = 2$ gives, with $\varphi$ the angle between the void edges
and the grid axes, for $0 \le \varphi \le 45^\circ$ and by symmetry at every angle, $c_\perp \ge (2 - \tfrac12\sin 2\varphi)(\sin\varphi +
\cos\varphi)\,h \ge 2h = m$, with equality for edges parallel to the grid,
so such an edge at $d = m$ is the worst straight edge. The
\corlOcsLandRmsIdm\,cm is the RMS of $\|\mathbf e\|$, so $\sigma =
\sqrt2$\,cm, and at $d = \corlMarginCm$\,cm, the margin, this probability
is $\Phi(-\sqrt2) \approx 0.079$. If ten of the roughly 97 footholds of a
sweep trial (\corlOcsStepsPerMFull{} steps per metre over
\corlOcsTrackM\,m) were this close to an edge, with independent errors,
the trial would stay void-free with probability only $(1 - 0.079)^{10}
\approx 0.44$. At a pillar corner, with voids beyond two perpendicular
edges at distance $d$, a foothold lands in a void with probability $1 - (1
- 0.079)^2 \approx 0.15$. We have
not measured how close GaitNet's footholds come to voids.
\end{remark}

\subsection{Clearance of the Grown Mask}

\looseness=-1 Model the terrain rays as samples on the lattice $\mathbf
r_0 + h\mathbb Z^2$ in the yaw frame of the decision, a sample being unsafe
if it lies in the closed void set $\mathcal V$ (exact perception and
registration). A candidate is a sample with no unsafe sample within
$\infty$-norm distance $kh$; this is the grown mask, a $(2k{+}1)\times
(2k{+}1)$ max-pool without padding, so every lattice point within $kh$ of a
candidate is sampled. Let
\[
  \rho_{\mathcal V} := \sup_{\mathbf p \in \mathcal V}\;
  \min\big\{\|\mathbf p - \mathbf r\|_\infty :
  \mathbf r \in (\mathbf r_0 + h\mathbb Z^2) \cap \mathcal V\big\}
\]
be the largest distance from a void point to the nearest lattice point in
the void (with $\rho_{\mathcal V} = \infty$, and Lemma~\ref{lem:idm-grid}(a)
empty, if part of the void holds no lattice point).

\begin{lemma}[Clearance]
\label{lem:idm-grid}
(a)~Every candidate $\mathbf u^\star$ has $d(\mathbf u^\star) \ge (k+1)h -
\rho_{\mathcal V}$.
(b)~If $\mathcal V$ is a union of closed rectangles with sides parallel to
the grid axes and at least $h$ long, then $\rho_{\mathcal V} \le h$, so
$d(\mathbf u^\star) \ge kh = m$.
(c)~If $\mathcal V$ is a union of closed squares of side at least
$2\sqrt2\,h$ whose edges make an angle $\varphi$ with the grid axes, then
$\rho_{\mathcal V} \le (1 + \tfrac12|\sin 2\varphi|)\,h$, so $d(\mathbf
u^\star) \ge (k - \tfrac12|\sin 2\varphi|)\,h$. For squares at different
angles, take the largest $|\sin 2\varphi|$.
\end{lemma}

\begin{idmproof}{Lemma~\ref{lem:idm-grid}}
(a)~Take a void point $\mathbf p$ and a lattice point $\mathbf r$ in the
void with $\|\mathbf p - \mathbf r\|_\infty \le \rho_{\mathcal V}$ (a
nearest one exists, as only finitely many lattice points lie within any
distance). If $\|\mathbf r - \mathbf u^\star\|_\infty \le kh$, then
$\mathbf r$ would be sampled, hence unsafe, within $kh$ of the candidate
$\mathbf u^\star$, which the mask excludes. So $\|\mathbf r - \mathbf
u^\star\|_\infty > kh$, and since both are lattice points, $\|\mathbf r -
\mathbf u^\star\|_\infty \ge (k+1)h$. By the triangle inequality, $\|\mathbf
p - \mathbf u^\star\|_\infty \ge (k+1)h - \rho_{\mathcal V}$ for every void
point $\mathbf p$.

(b)~Work on one axis. Let $\mathbf p$ lie in a rectangle whose side along
that axis is $[a,b]$, $b - a \ge h$, and let $t^- \le p \le t^+$ be the
lattice coordinates next to $p$, so $t^+ - t^- \le h$. If $t^- \ge a$, then
$t^- \in [a,b]$. Otherwise $t^+ \le t^- + h < a + h \le b$ and $t^+ \ge p \ge
a$, so $t^+ \in [a,b]$. Either way, a lattice coordinate in $[a,b]$ lies
within $h$ of $p$. Doing this on both axes gives a lattice point in the
rectangle within $\infty$-norm distance $h$ of $\mathbf p$.

(c)~\emph{Step 1: reduction.} Take $\mathbf p$ in a square $\mathcal Q
\subseteq \mathcal V$ of side $s \ge 2\sqrt2\,h$. Reflections in the grid
axes and the diagonal map the lattice to a translate of itself and
preserve $\|\cdot\|_\infty$, so we may take $h = 1$ and $0 \le \varphi \le
45^\circ$, and, moving the lattice, $\mathbf p = \mathbf 0$. Let $\mathbf d_1
= (\cos\varphi, \sin\varphi)$ and $\mathbf d_2 = (-\sin\varphi,
\cos\varphi)$ be the edge directions. $\mathcal Q$ extends at least $s/2$
from $\mathbf p$ along one of $\pm\mathbf d_1$ and one of $\pm\mathbf d_2$.
Rotations by multiples of $90^\circ$ about $\mathbf p$ permute these four
signed directions cyclically, keep $\varphi$, and also map the lattice to
a translate of itself, so we may take the two directions to be $+\mathbf
d_1$ and $+\mathbf d_2$. Then $\mathcal Q$ contains the quarter square
$\mathcal W := \{\gamma_1 \mathbf d_1 + \gamma_2\mathbf d_2 : 0 \le
\gamma_1,\gamma_2 \le s/2\}$.

\emph{Step 2: an axis-aligned unit square inside $\mathcal W$.} Let $B :=
[x_0, x_0+1] \times [y_0, y_0+1]$ with $x_0 = -\sin^2\varphi$ and $y_0 =
\sin\varphi\cos\varphi$. Its lower corners are $(\cos^2\varphi, y_0) =
\cos\varphi\,\mathbf d_1$ and $(x_0, y_0) = \sin\varphi\,\mathbf d_2$, and
its upper corners add $(0,1) = \sin\varphi\,\mathbf d_1 +
\cos\varphi\,\mathbf d_2$. In the coordinates $(\gamma_1,\gamma_2)$, the
four corners are $(\cos\varphi, 0)$, $(0, \sin\varphi)$, $(\cos\varphi +
\sin\varphi, \cos\varphi)$ and $(\sin\varphi, \sin\varphi + \cos\varphi)$:
all have $0 \le \gamma_1, \gamma_2 \le \sin\varphi + \cos\varphi \le \sqrt2
\le s/2$. By convexity, $B \subseteq \mathcal W \subseteq \mathcal Q$, and
every point of $B$ lies within $\infty$-norm distance $\max(\sin^2\varphi,
\cos^2\varphi, 1 + y_0) = 1 + \tfrac12\sin 2\varphi$ of $\mathbf 0$.

\emph{Step 3: a lattice point in $B$.} A closed axis-aligned unit square
contains a point of every translate of $\mathbb Z^2$, because a closed unit
interval contains a point of every translate of $\mathbb Z$. So $\mathcal Q$
contains a lattice point within $\infty$-norm distance $1 + \tfrac12\sin
2\varphi$ of $\mathbf p$. As $\mathbf p$ was any void point, this bounds
$\rho_{\mathcal V}$, and part (a) gives the clearance.
\end{idmproof}

\looseness=-1 Our tiles, $12.5h$ wide in Isaac and Gazebo, meet the side
condition $2\sqrt2\,h \approx 2.8h$. With $k = 2$ and $h = 1$\,cm, every
candidate therefore has a clearance of at least $(2 - \tfrac12|\sin
2\varphi|)$\,cm, that is, $2$\,cm for edges parallel to the grid and
$1.5$\,cm at $45^\circ$. The lemma assumes exact terrain. It ignores the
foot's size, and at deployment, where the grid is resampled from a
depth-camera map that keeps its own distance from unwalkable cells
(App.~\ref{app:deployment}), it does not cover that map's errors.
